\section{Setup and notation}
\label{sec:setup}

\subsection{Frames}

We work entirely in the plane. The world frame is inertial with axes
$(\hat{x}_w, \hat{y}_w)$. The body frame is rigidly attached to the vehicle
at the geometric centre of the wheel axle, with $\xbody$ pointing forward
(the direction of travel when both wheels rotate equally) and $\ybody$
pointing to the left:
\begin{equation}
    \xbody = \begin{pmatrix}\cos\theta \\ \sin\theta\end{pmatrix},
    \qquad
    \ybody = \begin{pmatrix}-\sin\theta \\ \phantom{-}\cos\theta\end{pmatrix}.
    \label{eq:body-axes}
\end{equation}

\subsection{Configuration and velocities}

The configuration is
\begin{equation}
    q \;=\; (p_x,\, p_y,\, \theta) \;\in\; \R^{2} \times \Sone,
\end{equation}
where $(p_x, p_y)$ is the world-frame position of the wheel-axle midpoint
and $\theta$ is the heading. We use the body twist $(v, \omega)$ as the
dynamical coordinates for translational and angular velocity:
\begin{align}
    v       &\;=\; \dot p_x \cos\theta + \dot p_y \sin\theta, \\
    \omega   &\;=\; \dot\theta.
\end{align}
Once the no-side-slip constraint of \S\ref{sec:constraints} is imposed,
this reduces to $(\dot p_x, \dot p_y) = v\,\xbody$.

\subsection{Geometric and actuator parameters}

\begin{description}
    \item[Track width $\trackw$] --- distance between the left and right
        wheel contacts. Default \(0.2\,\mathrm{m}\).
    \item[Wheel radius $\rwheel$] --- used only when converting between
        commanded wheel angular velocity $\dot\phi$ and rim velocity
        $v_w = \rwheel\,\dot\phi$.
    \item[Motor time constant $\taum$] --- first-order lag from commanded
        to realised body velocities. Default \(0.05\,\mathrm{s}\).
\end{description}

\subsection{Wheel inputs}

The control input is the pair of commanded wheel \emph{rim} velocities
$u = (v_L,\, v_R)$, with $v_L = \rwheel\,\dot\phi_L$ and similarly for the
right wheel. We use $v_L,\, v_R$ rather than $\dot\phi_L,\, \dot\phi_R$
throughout because the wheel radius is absorbed at the controller boundary
and never appears in the body dynamics. The \texttt{Control<Scalar>} type
in \texttt{types.hpp} stores exactly this pair.
